\subsection{The ordered cover of all intermediate arcs}

Saturation supplies the intermediate paths. We must now order their
local descriptions so that separate endpoint bounds become a bound
on one interval. The whole-passage tests will give a regular prefix
followed by at most one terminal union. At the parabolic corner this
requires two connections: from the finite family to the wedge, and
from the wedge to the secondary family.

\begin{lemma}[The two parabolic connections]
\label{rw:fibre:parabolic-connections}
Near $\mathcal P$ the finite regular family connects by an actual
open overlap to the wedge, possibly through a constant-sign lower
loss interval. The wedge in turn has an actual open overlap with
the regular secondary family. The construction applies to the
nonperiodic returning arcs in Proposition~\ref{rw:saturation}.
\end{lemma}
\begin{proof}
Consider first the incoming lower double boundary
$\bar a_*=2\sqrt{\bar B_*}$ and put
$r=1/\sqrt{\bar B_*}$. The limiting upper return has $q=p$.
For $0<p<r$,
\[
 A_-(p)-A_+(p)
 =2\bar a_*\int_0^p
               \frac{X^2}{F_+(X)F_-(X)}\,\dd X>0.
\]
Thus $K<0$ on every fixed compact regular initial interval,
uniformly in a smaller off-center neighborhood.
Let $C_w=\sqrt{1+2w_0}/w_0$. The limiting upper parabola
$Y=(p^2-X^2)/2$ meets $w=w_0$ at $H=C_w/p$.
Choose $L_1>2r$ so large that $C_w/L_1<H_0/2$.
The upper arcs for $L_1\le p\le2L_1$ form a compact regular
family; their transverse intersections persist and belong to
the wedge after the coefficients are reduced.

The incoming action tends to infinity as its first double root
is approached from below. Choose $R<r$ with
\[
 A_-(R)>\max_{p\le2L_1}A_+(p)+3,
\]
where the maximum is on the retained regular outgoing interval.
Then choose a fixed $p_1<r$ with $p_1>R+2\zeta$ for some
$\zeta>0$. Uniform convergence of the regular upper map to
$q=p$ on $[p_1,2L_1]$ gives $q(p)>R+\zeta$ there.
Lemma~\ref{rw:finite:incoming-contraction} now gives $K<0$
on its entire physical passing part, even beyond the lower
root where the ordinary regular action formula no longer applies.
This interval overlaps both the regular initial family and the
wedge. The corresponding lower slow action in the wedge has
the same strict sign and dominates its bounded remainder.
Reflection gives the other lower double boundary and reverses
the sign. Away from these boundary neighborhoods the lower
factors have a common positive margin on every retained finite
interval, so the regular family connects directly to the wedge.

For the wedge--secondary overlap, choose the fixed $L$ in secondary
routing large enough before decreasing any scale. Require
$2/L<\delta$ and $4/L^2<\delta$. At
$w=w_0$, $H=k_0\lambda L$, with $1/2<k_0<1$, the exact
secondary coordinates are
\[
 \widehat x=\frac1{k_0Lw_0},\qquad
 \widehat z=\frac1{k_0^2L^2w_0}.
\]
The auxiliary linear flow defining the actual tags gives
$P,Q=O(C_w/L)$ uniformly for bounded $(C,M)$.
At $C=M=0$ the values are $C_w/(k_0L)$; the perturbation
on these short arcs is of higher order in $1/L$.
Increase $L$ so all these tags lie in one compact regular
interval below the finite secondary thresholds.
The wedge inequalities are strict, and the same actual arcs
meet the secondary sections before reaching $w=w_0$.
Transversality persists as $k_0$ varies on a smaller open subinterval
of $(1/2,1)$. Thus an open family of the same returning arcs has both
descriptions, including the nonclosing comparisons.
\end{proof}

\begin{proposition}[An ordered cover on each physical fibre]
\label{rw:ordered-cover}
After a common original parameter shrink, every physical interval
of Proposition~\ref{rw:saturation} is covered by a connected
regular initial part and an ordered union of the local passing
charts of Sections~7--10. Their adjacent members overlap on
actual orbits. At a joint repeated center one contact member and
the full separated member form a single union. Near $\mathcal P$
the initial part includes the finite connection, the wedge and
the regular secondary part.
After grouping overlapping descriptions of the same end, there
is at most one terminal union after this initial part.
It may contain a constant-sign lower-loss connection followed
by a singular upper end; that connection is retained as its
own piece.

The coefficient neighborhoods and all cutoffs can be selected
from a finite family, independently of the chosen outer cycle.
The order includes signed split-root faces and the repeated
one-infinite interval. Each local derivative estimate is applied
on its stated whole-passage domain. This proposition makes no
sum of local zero bounds.
\end{proposition}
\begin{proof}
Fix the finite derivative orders of Sections~7--10 and retain their
larger construction domains before reducing sections or coefficient
neighborhoods. Saturation supplies actual upper and lower returns
with positive central heights on $[p_0,p_\gamma]$.
Since $q'>0$, their endpoints satisfy $p\le p_\gamma$,
$q\le q_\gamma$ and inherit the outer lower cones.

The key ordering fact is that a lost passing component cannot reappear.
The exact lower factor on the fixed phase interval is
\[
 F_I(x)=\bar B-\bar ax+(1+\epsilon\bar m)x^2
 -\epsilon\nu(1+2\epsilon\bar m)x^3H(x)
 +\epsilon^3\nu^2\bar m x^4H(x)^2.
\]
By Proposition~\ref{rw:phase}, it is strictly convex and positive
near the fixed exit $x_R$. Lemma~\ref{rw:transitions:lower-sections}
requires $F_I(s)>0$ at every point between the actual entry and
$x_R$. Thus only the component adjoining $x_R$ passes: if a root
is present, the entry must lie to the right of its last root.
The entry decreases as its actual endpoint increases, since
$T=1/p$ or $1/q$ and $t_T>0$. Its allowed endpoints consequently
form one outward initial interval. Entry positivity in the other
component would not suffice. Also $F_I'(0)=-\bar a$ and its
reflected derivative is $\bar a$; strict convexity permits an
interior positive double minimum on at most one of the two sides.

The primary upper slow tests have the same one-sided form. For a repeated
factor $v^2+\mu_U$, passage to the positive fixed exit requires
$\mu_U>0$ or $v>\sqrt{-\mu_U}$; its actual finite entry decreases
with $p$. This retains the negative entries beyond the repeated
threshold whenever they pass. For the weak factor $u(\gamma+u)$,
the whole domain is $\max(0,-\gamma)<u<u_0$, and
$p=\hat p-t(u)$ with $t'>0$. At joint contact both calibrated
entries decrease, so the two whole-path tests are intersections
of one-sided restrictions. General horizontal factors use the
whole lower test above; a smaller monotone block alone permits
an endpoint test. Signed auxiliary coefficient comparisons do not
remove physical inputs, whose complete central heights are positive
by saturation.

The increasing upper graph and its reflected inverse are those of
Lemma~\ref{rw:finite:upper-return}. Its nonregular ends are classified
in Lemma~\ref{rw:fibre:candidate-space}, with the full repeated
interval $p\ge P$, $q=\infty$ retained. Following the graph outward
gives these terminal templates:
\begin{itemize}
\item At a regular finite upper end, the first lower loss is a
single root or the simultaneous simple-root pair. A simple root
terminates passage; a double retains its signed slow alternatives.
\item At a finite upper saddle, the lower paths are regular, in
one or both contact neighborhoods, or have one incoming double-loss
connection before the saddle. An outgoing double with regular
incoming outer data is excluded by
Lemma~\ref{rw:finite:action-order}, applied to the outer cycle.
The incoming connection is the whole constant-sign interval of
Lemma~\ref{rw:finite:incoming-contraction}.
\item At a one-infinite end, a lower root can precede, coincide
with, or follow the upper threshold. These give respectively the
lower connection, the contact neighborhood, or the upper-tail piece
followed by the lower loss. A root escaping to infinity retains
its finite-root and cubic descriptions in the same family.
\item At two infinite endpoints the lower cones exclude a prior
simple root. Section~8 supplies the strict-focus, one lower-double,
repeated, opposing-slow and cubic alternatives; Section~9 supplies
the joint contact interval and its full separated tail.
\end{itemize}
A finite saddle has one returning quadrant, while the strict,
repeated and weak preparations have the through components just
described. Together with the monotone entry tests, this excludes a
second independent upper end or a return after a split-root
obstruction. The pieces of each listed end therefore form one
terminal union, retaining a preceding constant-sign lower connection
as a separate piece within it.

Near $\mathcal P$, use the actual first exits and tags of
Proposition~\ref{rw:secondary:exhaustion}. Routing confines a
returning arc to the fixed secondary rectangle and uniqueness
excludes another middle excursion. The normalized face
$\max(|C|,\sqrt{|M|})=1$ has the finite-saddle, strict real,
repeated, weak and focus alternatives, including $1/P=0$ and
$1/Q=0$. The repeated row retains its full label interval.
At its weak end the actual tag is $P=P_s-\psi(t(u))$ with
$\psi'>c$; again outward motion decreases $u$.
The preceding lemma inserts the finite connection, wedge and
regular secondary prefix. The same first-threshold reasoning then
gives one secondary terminal family; a lower-double overlay can
add its retained connection but cannot add an opposing double or
another returning quadrant.

Retain actual open seams before reducing the scale: a fixed positive
deficit crosscut between finite regular and singular pieces; the
joint contact/separated overlap $P+d<p<P+2d$ with both gaps
positive; and a fixed reciprocal annulus before either slow block,
using Lemma~\ref{rw:transitions:common-input}. These connectors
are transverse with bounded inverses. The finite lower-double
bridge has the strict action margin of the preceding lemma.
That lemma also supplies the finite--wedge and wedge--secondary
transverse overlaps. Within the secondary family retain its fixed
deficit crosscuts and reciprocal annuli. Adjacent strict, repeated
and weak faces use their own uniform preparations; no estimate at
a fixed positive eigenvalue ratio is extended across zero or infinity.
All seams compare the same original central function.

Apply Lemma~\ref{rw:fibre:compact-selection} to the joint compactification
of Lemma~\ref{rw:fibre:candidate-space}: the intermediate copy has
independent bounds $p\ge p_0$, $q\ge q_0>0$, its own first-exit tags,
and $p\le p^{\rm out}$, $q\le q^{\rm out}$. Each compact limiting
fibre is covered by the ordered pieces and their whole regular prefix.
The retained crosscuts, action margins and transverse seams make this
an open cover, including reciprocal radial ends; the one-sided entry
tests preserve order. Positive endpoint limits from the initial energy
classes use the regular enlargement in
Lemma~\ref{rw:fibre:candidate-cover}. Thus the lemma gives common
cutoffs after all sections, derivative budgets and overlap margins
have been fixed.
\end{proof}
